\section{方法与匹配路径对照}
\label{sec:method}

\subsection{任务与一层时空上下文}
设窗口$X\in\mathbb R_{\geq0}^{T\times F}$，$T=50$，
掩码$M_{ti}=1$表示可见。任务是\emph{离线整窗补全}，允许读取前后两侧观测。
完整拟合段确定全局均值$\mu$和总体标准差$\sigma$（下限$10^{-8}$），
以及每流自身和相关系数最高的至多八个正Pearson相关邻流的集合。
全局方差由流内方差与流均值间方差合成，不采用逐流标准化。
当前窗口的隐藏值不进入编码，标准化输入为
\begin{equation}
 z_{ti}=\begin{cases}(X_{ti}-\mu)/\sigma,&M_{ti}=1,\\0,&M_{ti}=0.\end{cases}
 \label{eq:observed-input}
\end{equation}
实现先按掩码清除隐藏输入，不依赖该位置的真值。

槽位$t/49$经可学习映射，与流身份嵌入及正弦位置编码构成32维$p_{ti}$。
可见位置加入数值编码，经归一化和数值跳连后，
由一层32维SPIN时空块得到$h_{ti}$\cite{marisca2022spin}。
本流时间注意力屏蔽同刻对角项；空间注意力排除自身流，
逐邻流在源时间维归一化后聚合。所有查询位置均保留，只有真实可见位置作为注意力源；
空支持读出为零，并保留根节点跳连。
由$h_{ti}$经带偏置仿射投影获得128维查询，分为四个32维块；
每块$v$按$v/\max(\|v\|_2,10^{-6})$归一化，拼接记为$q(h_{ti})$。
整窗上下文和双向读取均不提供在线因果保证。

\subsection{双向本流观测直读}
对目标$(t,i)$，从$s\leq t$与$s\geq t$的本流可见位置中，
分别选距离最近的至多两个，记为$\mathcal S^{\rightarrow}_{ti}$和$\mathcal S^{\leftarrow}_{ti}$。
可见目标自身可进入两侧集合；待补全位置不会读取自身隐藏值。
两方向共享带偏置的$32\to16$仿射投影$\phi_Q^D,\phi_K^D$，评分为
\begin{equation}
 \ell_{tis}=\frac{\phi_Q^D(p_{ti})^\top\phi_K^D(p_{si})}{\sqrt{16}}
             -\log(1+|t-s|).
 \label{eq:direct-score}
\end{equation}
在每个非空方向集合内作softmax，得到$\alpha^{r}_{tis}$，其中$r\in\{\rightarrow,\leftarrow\}$。
每方向仅提取三个统计量：
\begin{equation}
 \begin{aligned}
 v^r_{ti}&=\sum_{s\in\mathcal S^r_{ti}}\alpha^r_{tis}z_{si},\\
 \delta^r_{ti}&=\sum_{s\in\mathcal S^r_{ti}}\alpha^r_{tis}|t-s|/49,\\
 a^r_{ti}&=\mathbb I[\mathcal S^r_{ti}\ne\varnothing],\qquad
 d^r_{ti}=A[v^r_{ti};\delta^r_{ti};a^r_{ti}],
 \end{aligned}
 \label{eq:direct-statistics}
\end{equation}
其中两方向共享的$A\in\mathbb R^{128\times3}$无偏置，空支持时三个统计量及读出均为零。
因此$\operatorname{rank}(A)\leq3$：128维是统计量的嵌入宽度，
不意味着提取128个独立的局部特征。
给定参数后，权重由位置、距离和可见支持决定，不由当前流量数值决定；
数值经加权、投影和解码后输出，不能称为保证保幅的原值复制。

\subsection{解码与训练目标}
组合模型将两向Direct读出和上下文查询拼接，
由$384\to128\to1$的GELU解码器$D$输出原量纲预测：
\begin{equation}
 R_{ti}=\mu+\sigma D([d^{\rightarrow}_{ti};d^{\leftarrow}_{ti};q(h_{ti})]).
 \label{eq:decoder}
\end{equation}
推理时对缺失位置作非负截断，并回填可见真值。
所有保留模块联合训练，解码末层零初始化。
对批次缺失位置集合$\Omega$，使用未截断的原量纲L1损失
\begin{equation}
 \mathcal L=\frac{\sum_{(b,t,i)\in\Omega}|R_{bti}-X_{bti}|}{C|\Omega|},
 \quad C=\operatorname{mean}|X_{\rm fit}|.
 \label{eq:missing-loss}
\end{equation}
观测回填和非负截断均不参与该损失。

\subsection{两个单路径对照及实现边界}
\textbf{Context-only}删除Direct的评分与投影参数，
保留SPIN、$q(h)$和原解码器，输入变为$[0;0;q(h)]$。
\textbf{Direct-only}删除数值上下文编码器与SPIN聚合，
保留位置编码、Direct与解码器，并将查询改为$q(p)$，即输入$[d^{\rightarrow};d^{\leftarrow};q(p)]$。
后者的目标流预测不接收其他流当前窗口的数值、掩码或图信息，
但仍沿用共同的拟合期全局标准化。
两者均独立训练，不是推理时临时关闭分支。

同一种子先按组合模型的完整顺序初始化，再删除停用模块，
使共有参数初值逐位一致。
这些对照检验整条信息路径的使用价值，\emph{不是等参数量或等有效容量比较}。
Context-only首层有$256\times128=32{,}768$个权重对应恒零输入，
数据损失对其梯度为零，虽仍受AdamW权重衰减，不能表达额外信息。
Context-only也不是官方SPIN的复现替代品。

实现方面，Direct先计算同流完整$T\times T$评分再掩蔽筛选，
因此逻辑上读取至多两个观测不构成时间线性复杂度保证。
SPIN也沿用密集时间配对及内部节点分块；本文只报告当前实现的实测成本。
历史Full另将两个方向的同步记忆读出分别与$d^r$相加，
其120轮结果仅用于展示额外记忆的取舍，不与本轮160轮路径对照混作同预算结论。
